% ch-reproducibility.tex -- numerical method: what bit-level reproducibility depends on
\chapter{Bit-level reproducibility}
\label{ch:reproducibility}

\textbf{This chapter is for developers who intend to modify the solver.} If you
only run Yuzen, read \cref{sec:linsolve:choosing} and stop there.

Yuzen is written so that \textbf{changing the number of OpenMP threads does not
change the result by a single bit}. The one exception is the \val{"amgcl"}
backend, whose internal reductions are non-deterministic
(\cref{sec:linsolve:choosing}). The property is itself one of the claims the
test suite makes: \code{ctest} registers three entries that run the same binary
in three separate processes --- once with the environment's own thread count,
once with \env{OMP\_NUM\_THREADS=1} and once with \env{OMP\_NUM\_THREADS=4}.

Four things in the code carry that guarantee. Every one of them looks like
something a tidy-minded reader would want to rewrite, and every one of them
breaks bit-identity if rewritten. They are set out one per section below, with
the remaining rules collected at the end.

If you are here because you are adding physics rather than auditing the code,
read this chapter first and then \cref{ch:extending}, which restates these four
as the rules a new term has to respect and gives the procedure for adding one.
\Cref{ch:layout} is the map of where the code discussed here lives.

\section{The deterministic reductions in \code{Vector}}
\label{sec:repro:reduction}

\code{\#pragma omp parallel for reduction(+:x)} decides how to split the
partial sums, and in what order to add them, from the number of threads.
Floating-point addition is not associative, so a different thread count gives
different low-order bits in an inner product or a norm. Those bits do not stay
small: they propagate into the iteration count of BiCGStab and into the
acceptance test of the time-step control, and surface as ``one thread and two
threads give different answers''.

The inner product (\code{operator,}), \code{Sum()} and \code{Norm()} of
\code{Vector} therefore split the array into \textbf{a fixed number of blocks
that does not depend on the thread count}, form a partial sum per block, and
add the blocks up in block order at the end. The number of blocks follows from
the array length alone --- at least 512 elements per block, at most 256 blocks
--- so the order of the additions is the same however many threads are running.

\begin{warn}[Do not change the order of summation]
  Do not replace these with a bare \code{reduction(+:...)}.
  \file{yuzen/tests/test\_vector.cpp} pins the behavior down.
\end{warn}

\code{Max}, \code{Min} and \code{AbsMax} need none of this: they are exactly
associative and exactly commutative, so a plain sequential loop is correct.

\section{The storage order of \code{CRS\_Matrix}}
\label{sec:repro:crs}

Within a row, the entries are stored \textbf{in insertion order}, and the first
entry of every row is always the diagonal. This is not a question of
presentation. It fixes the order of the floating-point additions in the
matrix--vector product, in \code{ScalingDiag} and in \code{Residual}.

\begin{warn}[Do not sort the entries of a row into ascending column order]
  Sorting changes how the sum rounds, and bit-identity is lost.
\end{warn}

Libraries that require ascending column indices within a row --- Eigen and
AMGCL --- are given a \textbf{sorted mirror}, kept separately from the matrix
proper. The permutation is built once, at the moment the sparsity structure is
finalized; after that each handover is an $O(\mathrm{nnz})$ gather of the
values.

\code{CRS\_Matrix} has two states: a construction phase, in which each row is a
separate variable-length array, so that rows are independent of one another and
insertion can safely be done under OpenMP; and a finalized phase holding flat
CRS. Any read-side call finalizes it automatically. One detail of the
construction phase is worth knowing: the boundary-condition flag is a
\code{char} rather than a \code{bool}, because \code{std::vector<bool>} is a
bit-packed proxy reference and is not thread-safe to write.

\section{The registration order of the physical terms}
\label{sec:repro:order}

\code{WellMixed\_HeightEquationBase::RegisterPhysics()} lists the physical
terms, and that list decides the order of the \code{+=} operations into the
matrix, which is the column insertion order of \cref{sec:repro:crs}. The order
is

\begin{plain}
LaplacePressure -> VerticalGravity -> ParallelGravity
                -> SolutoCapillarity -> Centrifugal -> DisjoiningPressure
\end{plain}

\begin{warn}[Do not reorder the list; add a new term at the end]
  Reordering changes the storage order, and with it the rounding of the sums
  described in \cref{sec:repro:crs}.
\end{warn}

The rule is worth stating positively as well, because it is what makes the
constraint tolerable: \textbf{as long as a new term is appended at the end, no
existing result changes by a single bit}. A term that is not enabled does not
even have its \code{SetMatrixA} called, so it contributes nothing to the order
of operations of any case that does not use it.

The velocity calculation \code{CalVelocity} is the one exception. Each term
writes into a velocity field of its own there, so the order in which they run
cannot affect the result.

\section{The two direction orderings}
\label{sec:repro:directions}

The kernels carry \textbf{two} orderings of the directions and use them in
different places \emph{deliberately} (\file{GridTraits.h}):

\begin{itemize}
  \item \code{T::Dir(d)} is axis order --- in 2-D, \code{Xm}, \code{Xp},
    \code{Ym}, \code{Yp}.
  \item \code{T::StencilDir(k)} is stencil order --- in 2-D, \code{Ym},
    \code{Xm}, \code{Xp}, \code{Yp}.
\end{itemize}

\code{SetMatrixA} sums each row in axis order. \code{SetJacobiMatrix} sums the
diagonal and the cross terms in stencil order. The diagonal of the capillary
pressure term contains both in a single expression:
\code{gBo*aB + gLo*aL + gRo*aR + gTo*aT} is in stencil order and
\code{sLo*bL + sRo*bR + sBo*bB + sTo*bT} is in axis order, on the two sides of
the same sum.

\begin{warn}[Follow the comment in the kernel about which ordering to use]
  Swapping one for the other breaks bit-identity. In particular, do not ``tidy
  up'' by making both sites use the same ordering.
\end{warn}

In 1-D the two orderings coincide, so the distinction is invisible; it shows up
only in 2-D. That is exactly what makes it easy to break without noticing.

\section{The remaining rules}
\label{sec:repro:other}

\begin{itemize}
  \item \textbf{A preconditioner's triangular substitutions must run
    sequentially.} Both the forward and the backward substitution carry a
    loop-carried dependency. Parallelizing them makes a thread read values
    another thread has not updated yet, which silently turns the
    preconditioner into a block Jacobi whose blocks are the thread
    decomposition --- so the preconditioner, and hence the iteration count,
    becomes thread-count dependent. The rule binds any preconditioner added to
    a backend of your own (\cref{ch:extending}). The block-Jacobi ILUT of
    \val{"eigen"} keeps it: its blocks follow from the size of the system,
    not from the threads, and each block's substitutions run sequentially
    (\cref{sec:linsolve:eigen}).
  \item \textbf{A decision taken inside a solve depends on numbers, not on
    threads.} Whether a light preconditioner's try is kept or dropped, the
    tolerance of a Newton pass and the start of the Newton iteration follow
    from the convergence of the solves and the fields
    (\cref{sec:linsolve:conc,sec:linsolve:ew,sec:time:newton}), all of them
    the same at any thread count with \val{"eigen"}.
  \item \textbf{\code{Vector} does not skip zero-initialization};
    \code{resizeFast()} does exactly what \code{resize()} does. Any path that
    can read an uninitialized value makes the result non-deterministic.
  \item \textbf{The compiler does not fuse multiplies and adds.} The build
    passes \code{-ffp-contract=off} to GCC and Clang. On a processor with a
    fused multiply-add (every AArch64 one, including Apple silicon), the
    compiler would otherwise turn \code{a*b + c} into one instruction where it
    sees fit, and may do so differently in two places that inline the same
    code. A run restarted from its state files then no longer matches the run
    made straight through. On x86-64 built without \code{-march} the flag
    changes nothing; do not remove it, and do not add an option that
    re-enables contraction (\code{-ffast-math}, \code{-ffp-contract=fast}).
  \item \textbf{Do not keep state in a function-local \code{static}.} State
    that ought to belong to an instance becomes state belonging to the
    translation unit: several cases can no longer be run in one process, and
    counters bleed between separate instances --- between the thickness
    equation and the concentration equation, for example.
\end{itemize}
